\section*{Bloque C. Casos con implicancias regulatorias y sector saneamiento}

\noindent\textit{Preguntas 66 a 90. Evalúan aplicaciones econométricas avanzadas al sector de agua y saneamiento: estimación de elasticidades de demanda, tarifas por bloques crecientes, fronteras estocásticas de costos (SFA) en EPS, diferencias en diferencias (Diff-in-Diff), regresión discontinua en SISFOH, WACC y modelos de datos de panel.}

\q{En la estimación econométrica de la demanda de agua potable residencial para fijar tarifas en SUNASS, la especificación comúnmente adoptada es de tipo log-log: $\ln(Q_i) = \alpha + \beta_P \ln(P_i) + \beta_Y \ln(Y_i) + \gamma Z_i + u_i$. Los parámetros $\beta_P$ y $\beta_Y$ representan directamente:}
\begin{opts}
    \item el consumo básico de subsistencia y el cargo fijo mensual.
    \item la elasticidad precio directa y la elasticidad ingreso de la demanda.
    \item la semielasticidad tarifaria media y el costo marginal.
    \item el margen operativo regulado y el factor de productividad.
\end{opts}

\q{En los sistemas de tarifas de agua por bloques crecientes (Increasing Block Tariffs - IBT) utilizados por las EPS peruanas, la estimación por MCO de la demanda de agua enfrenta un grave problema de endogeneidad porque:}
\begin{opts}
    \item el agua potable es un bien homogéneo cuya utilidad marginal decreciente impide estimar regresiones.
    \item el precio marginal y la cantidad consumida se determinan simultáneamente por el salto de bloque.
    \item la Superintendencia no publica los pliegos tarifarios vigentes de cada bloque en el diario El Peruano.
    \item el consumo de agua residencial es estacionalmente fijo e invariante ante cambios tarifarios mensuales.
\end{opts}

\q{Para corregir la endogeneidad de la tarifa marginal en modelos de demanda de agua con bloques crecientes, la literatura econométrica y los estudios regulatorios emplean como variable instrumental:}
\begin{opts}
    \item la tarifa del primer bloque o estructura tarifaria exógena de EPS comparables.
    \item la tasa de interés de referencia del BCRP y el encaje bancario comercial promedio.
    \item el índice de precios al consumidor de alimentos y bebidas del área metropolitana.
    \item el producto bruto interno per cápita nacional ajustado por paridad cambiaria.
\end{opts}

\q{SUNASS utiliza modelos de Análisis de Frontera Estocástica (Stochastic Frontier Analysis - SFA) para estimar los costos eficientes de las EPS. La ventaja econométrica clave de SFA frente a una regresión MCO estándar es que:}
\begin{opts}
    \item asume que todas las empresas prestadoras operan exactamente sobre la frontera con retorno del $100\%$.
    \item descompone el error en ineficiencia técnica no negativa (asimétrica) y ruido aleatorio simétrico.
    \item elimina la necesidad de especificar forma funcional de costos o calcular pruebas de significancia $t$.
    \item no requiere recopilar información contable ni estados financieros auditados de las empresas EPS.
\end{opts}

\q{Al estimar una función de costos translogarítmica multirregulada para las 50 EPS del Perú, se evalúan simultáneamente economías de escala y economías de densidad de clientes. La economía de densidad se define econométricamente como:}
\begin{opts}
    \item la tasa de depreciación de los micromedidores instalados por kilómetro.
    \item la variación porcentual del costo ante mayor volumen manteniendo fijos red y clientes.
    \item la reducción del costo medio cuando se fusionan dos EPS distantes.
    \item la relación técnica entre el gasto laboral y el costo de capital.
\end{opts}

\q{En una regresión de corte transversal de costos de operación (OPEX) que incluye a Sedapal junto a empresas provinciales pequeñas (como EPS Moyobamba o EPS Ilo), la presencia de heterocedasticidad es sumamente probable debido a que:}
\begin{opts}
    \item las empresas comparadas utilizan tecnologías distintas de micromedición ultrasónica.
    \item los residuos de la regresión siguen necesariamente una distribución uniforme discreta.
    \item todas las EPS aplican exactamente la misma estructura tarifaria para uso doméstico.
    \item la enorme diferencia de escala y dispersión entre Sedapal y las pequeñas EPS locales.
\end{opts}

\q{Para corregir la distorsión inferencial causada por la heterocedasticidad en la estimación de costos de EPS con muestras pequeñas a medianas ($n \approx 50$), la literatura econométrica recomienda utilizar la corrección de White modificada:}
\begin{opts}
    \item estimador sándwich asintótico HC0 sin ajuste de muestra.
    \item regresión de cuantiles no condicionada por mínimos desvíos.
    \item errores estándar convencionales de MCO homocedásticos.
    \item HC3 (Jackknife), que penaliza observaciones con alto leverage.
\end{opts}

\q{Al modelar series de tiempo mensuales de caudales de captación de agua en ríos de la vertiente del Pacífico para planificar el abastecimiento de EPS, la presencia de estacionalidad y persistencia intertemporal se modela mediante:}
\begin{opts}
    \item modelos SARIMA, que capturan autocorrelación serial y ciclos estacionales.
    \item análisis de componentes principales no lineales sobre series anuales.
    \item regresión lineal simple con intercepto nulo y tendencia determinística.
    \item modelos de selección muestral en dos etapas de corte transversal.
\end{opts}

\q{En modelos de series de tiempo de demanda de agua urbana donde se evalúa el impacto de episodios de El Niño, la aplicación de errores estándar HAC de Newey-West con $L$ rezagos permite:}
\begin{opts}
    \item garantizar que el estadístico de Durbin-Watson sea exactamente igual a 2.
    \item obtener inferencia robusta frente a correlación serial de choques y heterocedasticidad estacional.
    \item sustituir las series históricas pluviométricas de Senamhi mediante interpolación.
    \item eliminar por completo la tendencia poblacional y las variables ficticias.
\end{opts}

\q{En un panel de seguimiento de consumo de agua con micromedidores domiciliarios a lo largo de 5 años, el hecho de que los hogares con medidores defectuosos o manipulados sean retirados sistemáticamente de la base de datos genera un problema de:}
\begin{opts}
    \item violación del Teorema de Gauss-Markov por exceso de grados de libertad en la muestra.
    \item sesgo de atrición no aleatorio (Attrition Bias) correlacionado con el volumen consumido.
    \item multicolinealidad perfecta entre los efectos fijos temporales y la estructura tarifaria.
    \item discriminación de precios de segundo grado en la estimación de elasticidades cruzadas.
\end{opts}

\q{Para evaluar el impacto de un programa piloto de instalación masiva de micromedidores sobre el consumo de agua ejecutado por una EPS, la metodología cuasiexperimental estándar es:}
\begin{opts}
    \item aplicar la prueba de normalidad asintótica de Jarque-Bera sobre la distribución del consumo.
    \item comparar únicamente el promedio de consumo antes y después de la obra sin grupo de comparación.
    \item el estimador de Diferencias en Diferencias comparando grupos con y sin nuevos medidores.
    \item estimar una regresión lineal simple de corte transversal sin incluir término constante.
\end{opts}

\q{El supuesto de identificación indispensable para la validez causal del estimador de Diferencias en Diferencias en la evaluación del programa de micromedición es:}
\begin{opts}
    \item ausencia absoluta de correlación contemporánea entre las características de los hogares.
    \item que el coeficiente de determinación $R^2$ de la regresión auxiliar sea exactamente igual a 1.
    \item que la tarifa marginal media del agua potable iguale el costo medio de provisión del servicio.
    \item tendencias paralelas: sin intervención, tratados y controles habrían seguido igual trayectoria.
\end{opts}

\q{Para medir el impacto de la tarifa subsidiada del SISFOH sobre el bienestar de los hogares pobres, cuando el subsidio se asigna estrictamente si el puntaje de pobreza del hogar es menor a un umbral $c_0$, el método econométrico óptimo es:}
\begin{opts}
    \item Regresión Discontinua (RDD) nítida o difusa alrededor del umbral $c_0$.
    \item análisis de componentes principales de Hotelling sobre variables continuas.
    \item modelo de vectores autorregresivos (VAR) con restricciones de corto plazo.
    \item prueba de autocorrelación de Durbin-Watson para datos de corte transversal.
\end{opts}

\q{Al modelar encuestas de SUNASS sobre la decisión discreta de un hogar de pagar o no una tarifa adicional por agua con 24 horas continuas de servicio ($Y_i \in \{0, 1\}$), la estimación por el Modelo Lineal de Probabilidad (MLP) mediante MCO presenta la debilidad de que:}
\begin{opts}
    \item impide por completo incluir variables ficticias regionales o calcular correlaciones lineales.
    \item predice probabilidades fuera de $[0, 1]$ y sufre heterocedasticidad, prefiriéndose Logit/Probit.
    \item exige un tamaño de muestra infinito para poder estimar los parámetros mediante matrices.
    \item fuerza a que todos los coeficientes estimados converjan a cero cuando hay constante.
\end{opts}

\q{En un modelo de costos de inversión de las EPS que incluye simultáneamente como regresores: longitud de redes ($Km$), número de conexiones de agua ($Conex$) y volumen total producido ($m^3$), es muy probable encontrar Factores de Inflación de Varianza ($VIF > 10$) debido a:}
\begin{opts}
    \item la presencia de subsidios cruzados entre categorías comerciales e industriales.
    \item la correlación serial negativa entre las observaciones de distintas regiones naturales.
    \item la falta de normalidad asintótica de los residuos en la prueba de heterocedasticidad.
    \item la fuerte multicolinealidad estructural de variables que miden escala operativa.
\end{opts}

\q{Al modelar los costos de mantenimiento de redes de una EPS, si se incluye el porcentaje de Agua No Facturada ($ANF$) como regresor explicativo, existe endogeneidad por causalidad inversa porque:}
\begin{opts}
    \item mayor ANF eleva costos de reparación, pero menor mantenimiento preventivo dispara fugas y ANF.
    \item el porcentaje de agua no facturada es un parámetro constante fijado por ley.
    \item no existen registros históricos estandarizados de agua no facturada en las EPS.
    \item el indicador ANF es una variable cualitativa imposible de medir con precisión.
\end{opts}

\q{En encuestas de hogares donde se evalúa el gasto en agua potable, muchas familias periurbanas carecen de conexión a la red pública y no reportan tarifa de EPS. Modelar el gasto en agua solo para los hogares conectados genera un problema econométrico de:}
\begin{opts}
    \item violación del supuesto de homocedasticidad por exceso de grados de libertad.
    \item autocorrelación de orden 12 entre los residuos de consumo mensual domiciliario.
    \item multicolinealidad perfecta entre la variable dependiente y la constante.
    \item sesgo de selección muestral, corregible mediante el procedimiento de Heckman.
\end{opts}

\q{Al estimar la productividad de las 50 EPS mediante un panel de datos durante el periodo regulatorio 2020-2025 ($Y_{it} = \alpha_i + X_{it}\beta + u_{it}$), el modelo de Efectos Fijos (Fixed Effects) es preferible al de MCO agrupado (Pooled) porque:}
\begin{opts}
    \item no requiere descontar grados de libertad al estimar los parámetros de pendiente del modelo.
    \item elimina automáticamente toda heterocedasticidad y autocorrelación serial en los errores.
    \item controla la heterogeneidad individual invariante en el tiempo aun correlacionada con $X_{it}$.
    \item asume que todas las empresas prestadoras comparten exactamente el mismo intercepto individual.
\end{opts}

\q{La prueba de Hausman en modelos de datos de panel compara el estimador de Efectos Fijos ($\hat{\beta}_{FE}$) con el de Efectos Aleatorios ($\hat{\beta}_{RE}$). Si la hipótesis nula $H_0: Cov(\alpha_i, X_{it}) = 0$ es rechazada estadísticamente ($p < 0.05$):}
\begin{opts}
    \item Efectos Aleatorios es inconsistente por correlación con $\alpha_i$, debiendo usarse Efectos Fijos.
    \item ambos estimadores resultan algebraicamente idénticos al estimador MCO agrupado (Pooled).
    \item el estimador de Efectos Aleatorios sigue siendo consistente y supera en eficiencia a Efectos Fijos.
    \item el modelo de panel queda subidentificado y no puede estimarse por método alguno.
\end{opts}

\q{En el cálculo del Costo Promedio Ponderado de Capital (WACC regulatorio) que fija SUNASS para remunerar las inversiones de las EPS, el costo del capital accionario ($K_e$) se estima habitualmente mediante el modelo Capital Asset Pricing Model (CAPM): $K_e = R_f + \beta_L (R_m - R_f) + RP$. En esta ecuación, el coeficiente $\beta_L$ se estima econométricamente mediante:}
\begin{opts}
    \item el logaritmo natural de la razón entre volumen facturado y volumen producido por Sedapal.
    \item el cociente contable promedio entre el activo fijo neto y el pasivo corriente de la EPS.
    \item la prueba de Durbin-Watson aplicada sobre la serie mensual de tasa de interés interbancaria.
    \item regresión de retornos de acciones reguladas comparables sobre el retorno del mercado.
\end{opts}

\q{Dado que las EPS peruanas son empresas estatales no cotizadas en bolsa, SUNASS utiliza Betas de empresas comparables internacionales (proxy) desapalancadas y reapalancadas mediante la fórmula de Hamada: $\beta_L = \beta_U [1 + (1 - T)(D/E)]$. Si la Beta desapalancada promedio es $\beta_U = 0.50$, la tasa del impuesto a la renta es $T = 30\%$ y la estructura óptima de endeudamiento es $D/E = 0.60$, ¿cuál es la Beta apalancada regulada $\beta_L$?}
\begin{opts}
    \item $\beta_L = 1.00$
    \item $\beta_L = 0.85$
    \item $\beta_L = 0.50$
    \item $\beta_L = 0.71$
\end{opts}

\q{Para modelar econométricamente el número de roturas de tuberías o incidentes de atoros de alcantarillado por sector hidrográfico ($Y_i = 0, 1, 2, \dots$, datos de conteo), el modelo más apropiado que supera la limitación de la regresión lineal es:}
\begin{opts}
    \item el modelo de paseo aleatorio con deriva y raíz unitaria estacional.
    \item la prueba de cointegración en dos etapas de Engle y Granger.
    \item regresión de Poisson o Binomial Negativo (si hay sobredispersión).
    \item el modelo de Mínimos Cuadrados Ordinarios simple sin intercepto.
\end{opts}

\q{En la detección forense de colusión en licitaciones públicas de obras de agua y saneamiento convocadas por el Estado peruano, los econometristas analizan anomalías estadísticas en las ofertas de los postores buscando:}
\begin{opts}
    \item que todas las empresas postoras coticen exactamente al costo marginal unitario.
    \item varianza anormalmente baja entre ofertas y alta correlación de precios entre postores.
    \item un coeficiente $R^2$ igual a cero en la ecuación del presupuesto referencial.
    \item una distribución estrictamente bimodal en los precios de mercado regionales.
\end{opts}

\q{Un regulador implementa un modelo de Nowcasting para proyectar en tiempo real la recaudación tarifaria de las EPS durante episodios de crisis macroeconómica. Este método econométrico se caracteriza por:}
\begin{opts}
    \item asume que el término de perturbación estocástica futuro es conocido sin margen de incertidumbre.
    \item proyecta variables estructurales a 50 años aplicando exclusivamente el Teorema de Coase.
    \item utiliza únicamente modelos estáticos de equilibrio general computable sin rezagos dinámicos.
    \item estima el estado contemporáneo usando datos de alta frecuencia previos al reporte contable.
\end{opts}

\q{En las auditorías regulatorias de calidad bacteriológica del agua potable realizadas por SUNASS y DIGESA, la determinación del tamaño mínimo de muestra de grifos fiscalizados ($n$) para contrastar cumplimiento normativo exige balancear:}
\begin{opts}
    \item la depreciación contable lineal de los equipos de laboratorio y los costos fijos administrativos.
    \item la tasa de descuento intertemporal del BCRP y la volatilidad del tipo de cambio spot.
    \item el coeficiente $R^2$ ajustado de la regresión y el factor de inflación de varianza máximo.
    \item el error Tipo I (falsa sanción a la EPS) y el error Tipo II (no detectar agua contaminada).
\end{opts}

